% ============================================================
% DAFTAR PUSTAKA — SCOPE: DIAS (BACKEND)
% Muhammad Dias Al-Izzat — D4 TRM PENS
% Proyek: Sketchbook Universe — CNN + K-Means + REST API
%
% Total: 29 referensi (20 General/Shared + 9 Dias-specific)
% Format: IEEE numbered, urut by appearance (Bab 1 -> Bab 3)
%
% Source: Sitasi_Final_v1.md (16 Juni 2026)
% Catatan:
%   - B02 Meng et al. MediaPipe TIDAK dimasukkan ke Daftar Pustaka Dias
%     karena MediaPipe tidak ada di scope Dias.
%
% Cara pakai:
%   % ============================================================
% DAFTAR PUSTAKA — SCOPE: DIAS (BACKEND)
% Muhammad Dias Al-Izzat — D4 TRM PENS
% Proyek: Sketchbook Universe — CNN + K-Means + REST API
%
% Total: 29 referensi (20 General/Shared + 9 Dias-specific)
% Format: IEEE numbered, urut by appearance (Bab 1 -> Bab 3)
%
% Source: Sitasi_Final_v1.md (16 Juni 2026)
% Catatan:
%   - B02 Meng et al. MediaPipe TIDAK dimasukkan ke Daftar Pustaka Dias
%     karena MediaPipe tidak ada di scope Dias.
%
% Cara pakai:
%   % ============================================================
% DAFTAR PUSTAKA — SCOPE: DIAS (BACKEND)
% Muhammad Dias Al-Izzat — D4 TRM PENS
% Proyek: Sketchbook Universe — CNN + K-Means + REST API
%
% Total: 29 referensi (20 General/Shared + 9 Dias-specific)
% Format: IEEE numbered, urut by appearance (Bab 1 -> Bab 3)
%
% Source: Sitasi_Final_v1.md (16 Juni 2026)
% Catatan:
%   - B02 Meng et al. MediaPipe TIDAK dimasukkan ke Daftar Pustaka Dias
%     karena MediaPipe tidak ada di scope Dias.
%
% Cara pakai:
%   % ============================================================
% DAFTAR PUSTAKA — SCOPE: DIAS (BACKEND)
% Muhammad Dias Al-Izzat — D4 TRM PENS
% Proyek: Sketchbook Universe — CNN + K-Means + REST API
%
% Total: 29 referensi (20 General/Shared + 9 Dias-specific)
% Format: IEEE numbered, urut by appearance (Bab 1 -> Bab 3)
%
% Source: Sitasi_Final_v1.md (16 Juni 2026)
% Catatan:
%   - B02 Meng et al. MediaPipe TIDAK dimasukkan ke Daftar Pustaka Dias
%     karena MediaPipe tidak ada di scope Dias.
%
% Cara pakai:
%   \input{daftar_pustaka_dias_v1.tex}
%   atau copy bagian thebibliography ke proposal master.
% ============================================================

\chapter*{DAFTAR PUSTAKA}
\addcontentsline{toc}{chapter}{DAFTAR PUSTAKA}

\begin{thebibliography}{99}

\bibitem{ref1}
D. T. K. Ng, J. K. L. Leung, S. K. W. Chu, and M. S. Qiao, ``Conceptualizing AI literacy: An exploratory review,'' \textit{Computers and Education: Artificial Intelligence}, vol.~2, p.~100041, 2021.

\bibitem{ref2}
V. Casal-Otero \textit{et al.}, ``AI literacy in K-12: A systematic literature review,'' \textit{International Journal of STEM Education}, vol.~10, p.~29, 2023.

\bibitem{ref3}
H. Khosravi \textit{et al.}, ``Explainable Artificial Intelligence in education,'' \textit{Computers and Education: Artificial Intelligence}, vol.~3, p.~100074, 2022.

\bibitem{ref4}
E. Mosqueira-Rey, V. Alonso-R\'ios, and B. Rubio-R\'ios, ``Human-in-the-loop machine learning: A state of the art,'' \textit{Artificial Intelligence Review}, vol.~56, no.~4, pp.~3005--3054, 2023.

\bibitem{ref5}
B. Memarian and T. Doleck, ``Human-in-the-loop in artificial intelligence in education: A review and entity-relationship (ER) analysis,'' \textit{Computers in Human Behavior: Artificial Humans}, vol.~2, p.~100053, 2024.

\bibitem{ref6}
L. Xu \textit{et al.}, ``Deep learning for free-hand sketch: A survey,'' \textit{IEEE Transactions on Pattern Analysis and Machine Intelligence}, vol.~45, no.~1, pp.~285--312, 2023.

\bibitem{ref7}
G. Z. L. Goh, D. Ho, and Z. A. Abas, ``Front-end deep learning web apps development and deployment: A review,'' \textit{Applied Intelligence}, vol.~53, pp.~15923--15945, 2023.

\bibitem{ref8}
D. Smilkov \textit{et al.}, ``TensorFlow.js: Machine learning for the web and beyond,'' in \textit{Proc. 2nd SysML Conf.}, 2019, pp.~1--10.

\bibitem{ref9}
A. J. Karran, S. Demazure, and C. Hudelot, ``Designing for confidence: The impact of visualizing artificial intelligence decisions,'' \textit{Frontiers in Neuroscience}, vol.~16, 2022.

\bibitem{ref10}
P. Pansri \textit{et al.}, ``Understanding student learning behavior through K-Means clustering,'' \textit{Education Sciences}, vol.~14, no.~12, p.~1291, 2024.

\bibitem{ref11}
A. Alzahrani \textit{et al.}, ``Identifying weekly student engagement patterns via K-Means clustering,'' \textit{Electronics}, vol.~14, no.~15, p.~3018, 2025.

\bibitem{ref12}
M. Videnovik \textit{et al.}, ``Game-based learning in computer science education: A scoping literature review,'' \textit{International Journal of STEM Education}, vol.~10, p.~54, 2023.

\bibitem{ref13}
K. K. Chan, K. M. Wan, and R. B. King, ``Performance over enjoyment? Effect of game-based learning on learning outcome and flow experience,'' \textit{Frontiers in Education}, vol.~6, p.~660376, 2021.

\bibitem{ref14}
R. S. Pressman, \textit{Software Engineering: A Practitioner's Approach}, 8th~ed. New York: McGraw-Hill, 2014.

\bibitem{ref15}
I. Sommerville, \textit{Software Engineering}, 10th~ed. Pearson, 2015.

\bibitem{ref16}
Sugiyono, \textit{Metode Penelitian Kuantitatif, Kualitatif, dan R\&D}. Bandung: Alfabeta, 2020.

\bibitem{ref17}
J. W. Creswell, \textit{Research Design: Qualitative, Quantitative, and Mixed Methods Approaches}, 4th~ed. Sage Publications, 2014.

\bibitem{ref18}
K. Ishikawa, \textit{Guide to Quality Control}, 2nd~ed. Tokyo: Asian Productivity Organization, 1986.

\bibitem{ref19}
J. Brooke, ``SUS: A `Quick and Dirty' Usability Scale,'' in \textit{Usability Evaluation in Industry}, pp.~189--194, Taylor \& Francis, 1996.

\bibitem{ref20}
A. Bangor, P. T. Kortum, and J. T. Miller, ``An empirical evaluation of the System Usability Scale,'' \textit{International Journal of Human-Computer Interaction}, vol.~24, no.~6, pp.~574--594, 2008.

\bibitem{ref21}
M. Yin, J. Wortman Vaughan, and H. Wallach, ``Understanding the effect of accuracy on trust in machine learning models,'' in \textit{Proc. 2019 CHI Conf. Human Factors in Computing Systems}, 2019, pp.~1--12.

\bibitem{ref22}
C. Batanero and R. \'Alvarez-Arroyo, ``Teaching and learning of probability,'' \textit{ZDM -- Mathematics Education}, vol.~56, pp.~5--17, 2023.

\bibitem{ref23}
Z. Zhan, Y. Tong, X. Lan, and B. Zhong, ``A systematic literature review of game-based learning in artificial intelligence education,'' \textit{Interactive Learning Environments}, vol.~32, no.~4, pp.~1189--1212, 2024.

\bibitem{ref24}
Q. V. Liao, D. Gruen, and S. Miller, ``Questioning the AI: Informing design practices for explainable AI user experiences,'' in \textit{Proc. 2020 CHI Conf. Human Factors in Computing Systems}, 2020, pp.~1--15.

\bibitem{ref25}
Y. Feldman-Maggor, M. Cukurova, C. Kent, R. Luckin, and M. Baur-Wolfson, ``The impact of explainable AI on teachers' trust and acceptance of AI EdTech recommendations,'' \textit{International Journal of Artificial Intelligence in Education}, vol.~35, no.~1, pp.~1--25, 2025.

\bibitem{ref26}
V. T. Huynh, T. T. Nguyen, T. V. Nguyen, and M. T. Tran, ``MobileNet-SA: Lightweight CNN with self attention for sketch classification,'' in \textit{Pacific-Rim Symposium on Image and Video Technology (PSIVT) 2023}, Springer, 2023, pp.~115--129.

\bibitem{ref27}
L. Li, C. Zou, Y. Zheng, Q. Su, H. Fu, and C. L. Tai, ``Sketch-R2CNN: An RNN-rasterization-CNN architecture for vector sketch recognition,'' \textit{IEEE Transactions on Pattern Analysis and Machine Intelligence}, vol.~44, no.~9, pp.~4671--4685, 2020.

\bibitem{ref28}
Google Creative Lab, ``Quick, Draw! The Data,'' Google, 2017. [Online]. Available: https://quickdraw.withgoogle.com/data

\bibitem{ref29}
F. Liu, M. Lu, and W. Zhang, ``Comparative analysis of deep learning-based models in sketch recognition,'' in \textit{Proc. ICAIC 2024}, Atlantis Press, 2024.

\end{thebibliography}

%   atau copy bagian thebibliography ke proposal master.
% ============================================================

\chapter*{DAFTAR PUSTAKA}
\addcontentsline{toc}{chapter}{DAFTAR PUSTAKA}

\begin{thebibliography}{99}

\bibitem{ref1}
D. T. K. Ng, J. K. L. Leung, S. K. W. Chu, and M. S. Qiao, ``Conceptualizing AI literacy: An exploratory review,'' \textit{Computers and Education: Artificial Intelligence}, vol.~2, p.~100041, 2021.

\bibitem{ref2}
V. Casal-Otero \textit{et al.}, ``AI literacy in K-12: A systematic literature review,'' \textit{International Journal of STEM Education}, vol.~10, p.~29, 2023.

\bibitem{ref3}
H. Khosravi \textit{et al.}, ``Explainable Artificial Intelligence in education,'' \textit{Computers and Education: Artificial Intelligence}, vol.~3, p.~100074, 2022.

\bibitem{ref4}
E. Mosqueira-Rey, V. Alonso-R\'ios, and B. Rubio-R\'ios, ``Human-in-the-loop machine learning: A state of the art,'' \textit{Artificial Intelligence Review}, vol.~56, no.~4, pp.~3005--3054, 2023.

\bibitem{ref5}
B. Memarian and T. Doleck, ``Human-in-the-loop in artificial intelligence in education: A review and entity-relationship (ER) analysis,'' \textit{Computers in Human Behavior: Artificial Humans}, vol.~2, p.~100053, 2024.

\bibitem{ref6}
L. Xu \textit{et al.}, ``Deep learning for free-hand sketch: A survey,'' \textit{IEEE Transactions on Pattern Analysis and Machine Intelligence}, vol.~45, no.~1, pp.~285--312, 2023.

\bibitem{ref7}
G. Z. L. Goh, D. Ho, and Z. A. Abas, ``Front-end deep learning web apps development and deployment: A review,'' \textit{Applied Intelligence}, vol.~53, pp.~15923--15945, 2023.

\bibitem{ref8}
D. Smilkov \textit{et al.}, ``TensorFlow.js: Machine learning for the web and beyond,'' in \textit{Proc. 2nd SysML Conf.}, 2019, pp.~1--10.

\bibitem{ref9}
A. J. Karran, S. Demazure, and C. Hudelot, ``Designing for confidence: The impact of visualizing artificial intelligence decisions,'' \textit{Frontiers in Neuroscience}, vol.~16, 2022.

\bibitem{ref10}
P. Pansri \textit{et al.}, ``Understanding student learning behavior through K-Means clustering,'' \textit{Education Sciences}, vol.~14, no.~12, p.~1291, 2024.

\bibitem{ref11}
A. Alzahrani \textit{et al.}, ``Identifying weekly student engagement patterns via K-Means clustering,'' \textit{Electronics}, vol.~14, no.~15, p.~3018, 2025.

\bibitem{ref12}
M. Videnovik \textit{et al.}, ``Game-based learning in computer science education: A scoping literature review,'' \textit{International Journal of STEM Education}, vol.~10, p.~54, 2023.

\bibitem{ref13}
K. K. Chan, K. M. Wan, and R. B. King, ``Performance over enjoyment? Effect of game-based learning on learning outcome and flow experience,'' \textit{Frontiers in Education}, vol.~6, p.~660376, 2021.

\bibitem{ref14}
R. S. Pressman, \textit{Software Engineering: A Practitioner's Approach}, 8th~ed. New York: McGraw-Hill, 2014.

\bibitem{ref15}
I. Sommerville, \textit{Software Engineering}, 10th~ed. Pearson, 2015.

\bibitem{ref16}
Sugiyono, \textit{Metode Penelitian Kuantitatif, Kualitatif, dan R\&D}. Bandung: Alfabeta, 2020.

\bibitem{ref17}
J. W. Creswell, \textit{Research Design: Qualitative, Quantitative, and Mixed Methods Approaches}, 4th~ed. Sage Publications, 2014.

\bibitem{ref18}
K. Ishikawa, \textit{Guide to Quality Control}, 2nd~ed. Tokyo: Asian Productivity Organization, 1986.

\bibitem{ref19}
J. Brooke, ``SUS: A `Quick and Dirty' Usability Scale,'' in \textit{Usability Evaluation in Industry}, pp.~189--194, Taylor \& Francis, 1996.

\bibitem{ref20}
A. Bangor, P. T. Kortum, and J. T. Miller, ``An empirical evaluation of the System Usability Scale,'' \textit{International Journal of Human-Computer Interaction}, vol.~24, no.~6, pp.~574--594, 2008.

\bibitem{ref21}
M. Yin, J. Wortman Vaughan, and H. Wallach, ``Understanding the effect of accuracy on trust in machine learning models,'' in \textit{Proc. 2019 CHI Conf. Human Factors in Computing Systems}, 2019, pp.~1--12.

\bibitem{ref22}
C. Batanero and R. \'Alvarez-Arroyo, ``Teaching and learning of probability,'' \textit{ZDM -- Mathematics Education}, vol.~56, pp.~5--17, 2023.

\bibitem{ref23}
Z. Zhan, Y. Tong, X. Lan, and B. Zhong, ``A systematic literature review of game-based learning in artificial intelligence education,'' \textit{Interactive Learning Environments}, vol.~32, no.~4, pp.~1189--1212, 2024.

\bibitem{ref24}
Q. V. Liao, D. Gruen, and S. Miller, ``Questioning the AI: Informing design practices for explainable AI user experiences,'' in \textit{Proc. 2020 CHI Conf. Human Factors in Computing Systems}, 2020, pp.~1--15.

\bibitem{ref25}
Y. Feldman-Maggor, M. Cukurova, C. Kent, R. Luckin, and M. Baur-Wolfson, ``The impact of explainable AI on teachers' trust and acceptance of AI EdTech recommendations,'' \textit{International Journal of Artificial Intelligence in Education}, vol.~35, no.~1, pp.~1--25, 2025.

\bibitem{ref26}
V. T. Huynh, T. T. Nguyen, T. V. Nguyen, and M. T. Tran, ``MobileNet-SA: Lightweight CNN with self attention for sketch classification,'' in \textit{Pacific-Rim Symposium on Image and Video Technology (PSIVT) 2023}, Springer, 2023, pp.~115--129.

\bibitem{ref27}
L. Li, C. Zou, Y. Zheng, Q. Su, H. Fu, and C. L. Tai, ``Sketch-R2CNN: An RNN-rasterization-CNN architecture for vector sketch recognition,'' \textit{IEEE Transactions on Pattern Analysis and Machine Intelligence}, vol.~44, no.~9, pp.~4671--4685, 2020.

\bibitem{ref28}
Google Creative Lab, ``Quick, Draw! The Data,'' Google, 2017. [Online]. Available: https://quickdraw.withgoogle.com/data

\bibitem{ref29}
F. Liu, M. Lu, and W. Zhang, ``Comparative analysis of deep learning-based models in sketch recognition,'' in \textit{Proc. ICAIC 2024}, Atlantis Press, 2024.

\end{thebibliography}

%   atau copy bagian thebibliography ke proposal master.
% ============================================================

\chapter*{DAFTAR PUSTAKA}
\addcontentsline{toc}{chapter}{DAFTAR PUSTAKA}

\begin{thebibliography}{99}

\bibitem{ref1}
D. T. K. Ng, J. K. L. Leung, S. K. W. Chu, and M. S. Qiao, ``Conceptualizing AI literacy: An exploratory review,'' \textit{Computers and Education: Artificial Intelligence}, vol.~2, p.~100041, 2021.

\bibitem{ref2}
V. Casal-Otero \textit{et al.}, ``AI literacy in K-12: A systematic literature review,'' \textit{International Journal of STEM Education}, vol.~10, p.~29, 2023.

\bibitem{ref3}
H. Khosravi \textit{et al.}, ``Explainable Artificial Intelligence in education,'' \textit{Computers and Education: Artificial Intelligence}, vol.~3, p.~100074, 2022.

\bibitem{ref4}
E. Mosqueira-Rey, V. Alonso-R\'ios, and B. Rubio-R\'ios, ``Human-in-the-loop machine learning: A state of the art,'' \textit{Artificial Intelligence Review}, vol.~56, no.~4, pp.~3005--3054, 2023.

\bibitem{ref5}
B. Memarian and T. Doleck, ``Human-in-the-loop in artificial intelligence in education: A review and entity-relationship (ER) analysis,'' \textit{Computers in Human Behavior: Artificial Humans}, vol.~2, p.~100053, 2024.

\bibitem{ref6}
L. Xu \textit{et al.}, ``Deep learning for free-hand sketch: A survey,'' \textit{IEEE Transactions on Pattern Analysis and Machine Intelligence}, vol.~45, no.~1, pp.~285--312, 2023.

\bibitem{ref7}
G. Z. L. Goh, D. Ho, and Z. A. Abas, ``Front-end deep learning web apps development and deployment: A review,'' \textit{Applied Intelligence}, vol.~53, pp.~15923--15945, 2023.

\bibitem{ref8}
D. Smilkov \textit{et al.}, ``TensorFlow.js: Machine learning for the web and beyond,'' in \textit{Proc. 2nd SysML Conf.}, 2019, pp.~1--10.

\bibitem{ref9}
A. J. Karran, S. Demazure, and C. Hudelot, ``Designing for confidence: The impact of visualizing artificial intelligence decisions,'' \textit{Frontiers in Neuroscience}, vol.~16, 2022.

\bibitem{ref10}
P. Pansri \textit{et al.}, ``Understanding student learning behavior through K-Means clustering,'' \textit{Education Sciences}, vol.~14, no.~12, p.~1291, 2024.

\bibitem{ref11}
A. Alzahrani \textit{et al.}, ``Identifying weekly student engagement patterns via K-Means clustering,'' \textit{Electronics}, vol.~14, no.~15, p.~3018, 2025.

\bibitem{ref12}
M. Videnovik \textit{et al.}, ``Game-based learning in computer science education: A scoping literature review,'' \textit{International Journal of STEM Education}, vol.~10, p.~54, 2023.

\bibitem{ref13}
K. K. Chan, K. M. Wan, and R. B. King, ``Performance over enjoyment? Effect of game-based learning on learning outcome and flow experience,'' \textit{Frontiers in Education}, vol.~6, p.~660376, 2021.

\bibitem{ref14}
R. S. Pressman, \textit{Software Engineering: A Practitioner's Approach}, 8th~ed. New York: McGraw-Hill, 2014.

\bibitem{ref15}
I. Sommerville, \textit{Software Engineering}, 10th~ed. Pearson, 2015.

\bibitem{ref16}
Sugiyono, \textit{Metode Penelitian Kuantitatif, Kualitatif, dan R\&D}. Bandung: Alfabeta, 2020.

\bibitem{ref17}
J. W. Creswell, \textit{Research Design: Qualitative, Quantitative, and Mixed Methods Approaches}, 4th~ed. Sage Publications, 2014.

\bibitem{ref18}
K. Ishikawa, \textit{Guide to Quality Control}, 2nd~ed. Tokyo: Asian Productivity Organization, 1986.

\bibitem{ref19}
J. Brooke, ``SUS: A `Quick and Dirty' Usability Scale,'' in \textit{Usability Evaluation in Industry}, pp.~189--194, Taylor \& Francis, 1996.

\bibitem{ref20}
A. Bangor, P. T. Kortum, and J. T. Miller, ``An empirical evaluation of the System Usability Scale,'' \textit{International Journal of Human-Computer Interaction}, vol.~24, no.~6, pp.~574--594, 2008.

\bibitem{ref21}
M. Yin, J. Wortman Vaughan, and H. Wallach, ``Understanding the effect of accuracy on trust in machine learning models,'' in \textit{Proc. 2019 CHI Conf. Human Factors in Computing Systems}, 2019, pp.~1--12.

\bibitem{ref22}
C. Batanero and R. \'Alvarez-Arroyo, ``Teaching and learning of probability,'' \textit{ZDM -- Mathematics Education}, vol.~56, pp.~5--17, 2023.

\bibitem{ref23}
Z. Zhan, Y. Tong, X. Lan, and B. Zhong, ``A systematic literature review of game-based learning in artificial intelligence education,'' \textit{Interactive Learning Environments}, vol.~32, no.~4, pp.~1189--1212, 2024.

\bibitem{ref24}
Q. V. Liao, D. Gruen, and S. Miller, ``Questioning the AI: Informing design practices for explainable AI user experiences,'' in \textit{Proc. 2020 CHI Conf. Human Factors in Computing Systems}, 2020, pp.~1--15.

\bibitem{ref25}
Y. Feldman-Maggor, M. Cukurova, C. Kent, R. Luckin, and M. Baur-Wolfson, ``The impact of explainable AI on teachers' trust and acceptance of AI EdTech recommendations,'' \textit{International Journal of Artificial Intelligence in Education}, vol.~35, no.~1, pp.~1--25, 2025.

\bibitem{ref26}
V. T. Huynh, T. T. Nguyen, T. V. Nguyen, and M. T. Tran, ``MobileNet-SA: Lightweight CNN with self attention for sketch classification,'' in \textit{Pacific-Rim Symposium on Image and Video Technology (PSIVT) 2023}, Springer, 2023, pp.~115--129.

\bibitem{ref27}
L. Li, C. Zou, Y. Zheng, Q. Su, H. Fu, and C. L. Tai, ``Sketch-R2CNN: An RNN-rasterization-CNN architecture for vector sketch recognition,'' \textit{IEEE Transactions on Pattern Analysis and Machine Intelligence}, vol.~44, no.~9, pp.~4671--4685, 2020.

\bibitem{ref28}
Google Creative Lab, ``Quick, Draw! The Data,'' Google, 2017. [Online]. Available: https://quickdraw.withgoogle.com/data

\bibitem{ref29}
F. Liu, M. Lu, and W. Zhang, ``Comparative analysis of deep learning-based models in sketch recognition,'' in \textit{Proc. ICAIC 2024}, Atlantis Press, 2024.

\end{thebibliography}

%   atau copy bagian thebibliography ke proposal master.
% ============================================================

\chapter*{DAFTAR PUSTAKA}
\addcontentsline{toc}{chapter}{DAFTAR PUSTAKA}

\begin{thebibliography}{99}

\bibitem{ref1}
D. T. K. Ng, J. K. L. Leung, S. K. W. Chu, and M. S. Qiao, ``Conceptualizing AI literacy: An exploratory review,'' \textit{Computers and Education: Artificial Intelligence}, vol.~2, p.~100041, 2021.

\bibitem{ref2}
V. Casal-Otero \textit{et al.}, ``AI literacy in K-12: A systematic literature review,'' \textit{International Journal of STEM Education}, vol.~10, p.~29, 2023.

\bibitem{ref3}
H. Khosravi \textit{et al.}, ``Explainable Artificial Intelligence in education,'' \textit{Computers and Education: Artificial Intelligence}, vol.~3, p.~100074, 2022.

\bibitem{ref4}
E. Mosqueira-Rey, V. Alonso-R\'ios, and B. Rubio-R\'ios, ``Human-in-the-loop machine learning: A state of the art,'' \textit{Artificial Intelligence Review}, vol.~56, no.~4, pp.~3005--3054, 2023.

\bibitem{ref5}
B. Memarian and T. Doleck, ``Human-in-the-loop in artificial intelligence in education: A review and entity-relationship (ER) analysis,'' \textit{Computers in Human Behavior: Artificial Humans}, vol.~2, p.~100053, 2024.

\bibitem{ref6}
L. Xu \textit{et al.}, ``Deep learning for free-hand sketch: A survey,'' \textit{IEEE Transactions on Pattern Analysis and Machine Intelligence}, vol.~45, no.~1, pp.~285--312, 2023.

\bibitem{ref7}
G. Z. L. Goh, D. Ho, and Z. A. Abas, ``Front-end deep learning web apps development and deployment: A review,'' \textit{Applied Intelligence}, vol.~53, pp.~15923--15945, 2023.

\bibitem{ref8}
D. Smilkov \textit{et al.}, ``TensorFlow.js: Machine learning for the web and beyond,'' in \textit{Proc. 2nd SysML Conf.}, 2019, pp.~1--10.

\bibitem{ref9}
A. J. Karran, S. Demazure, and C. Hudelot, ``Designing for confidence: The impact of visualizing artificial intelligence decisions,'' \textit{Frontiers in Neuroscience}, vol.~16, 2022.

\bibitem{ref10}
P. Pansri \textit{et al.}, ``Understanding student learning behavior through K-Means clustering,'' \textit{Education Sciences}, vol.~14, no.~12, p.~1291, 2024.

\bibitem{ref11}
A. Alzahrani \textit{et al.}, ``Identifying weekly student engagement patterns via K-Means clustering,'' \textit{Electronics}, vol.~14, no.~15, p.~3018, 2025.

\bibitem{ref12}
M. Videnovik \textit{et al.}, ``Game-based learning in computer science education: A scoping literature review,'' \textit{International Journal of STEM Education}, vol.~10, p.~54, 2023.

\bibitem{ref13}
K. K. Chan, K. M. Wan, and R. B. King, ``Performance over enjoyment? Effect of game-based learning on learning outcome and flow experience,'' \textit{Frontiers in Education}, vol.~6, p.~660376, 2021.

\bibitem{ref14}
R. S. Pressman, \textit{Software Engineering: A Practitioner's Approach}, 8th~ed. New York: McGraw-Hill, 2014.

\bibitem{ref15}
I. Sommerville, \textit{Software Engineering}, 10th~ed. Pearson, 2015.

\bibitem{ref16}
Sugiyono, \textit{Metode Penelitian Kuantitatif, Kualitatif, dan R\&D}. Bandung: Alfabeta, 2020.

\bibitem{ref17}
J. W. Creswell, \textit{Research Design: Qualitative, Quantitative, and Mixed Methods Approaches}, 4th~ed. Sage Publications, 2014.

\bibitem{ref18}
K. Ishikawa, \textit{Guide to Quality Control}, 2nd~ed. Tokyo: Asian Productivity Organization, 1986.

\bibitem{ref19}
J. Brooke, ``SUS: A `Quick and Dirty' Usability Scale,'' in \textit{Usability Evaluation in Industry}, pp.~189--194, Taylor \& Francis, 1996.

\bibitem{ref20}
A. Bangor, P. T. Kortum, and J. T. Miller, ``An empirical evaluation of the System Usability Scale,'' \textit{International Journal of Human-Computer Interaction}, vol.~24, no.~6, pp.~574--594, 2008.

\bibitem{ref21}
M. Yin, J. Wortman Vaughan, and H. Wallach, ``Understanding the effect of accuracy on trust in machine learning models,'' in \textit{Proc. 2019 CHI Conf. Human Factors in Computing Systems}, 2019, pp.~1--12.

\bibitem{ref22}
C. Batanero and R. \'Alvarez-Arroyo, ``Teaching and learning of probability,'' \textit{ZDM -- Mathematics Education}, vol.~56, pp.~5--17, 2023.

\bibitem{ref23}
Z. Zhan, Y. Tong, X. Lan, and B. Zhong, ``A systematic literature review of game-based learning in artificial intelligence education,'' \textit{Interactive Learning Environments}, vol.~32, no.~4, pp.~1189--1212, 2024.

\bibitem{ref24}
Q. V. Liao, D. Gruen, and S. Miller, ``Questioning the AI: Informing design practices for explainable AI user experiences,'' in \textit{Proc. 2020 CHI Conf. Human Factors in Computing Systems}, 2020, pp.~1--15.

\bibitem{ref25}
Y. Feldman-Maggor, M. Cukurova, C. Kent, R. Luckin, and M. Baur-Wolfson, ``The impact of explainable AI on teachers' trust and acceptance of AI EdTech recommendations,'' \textit{International Journal of Artificial Intelligence in Education}, vol.~35, no.~1, pp.~1--25, 2025.

\bibitem{ref26}
V. T. Huynh, T. T. Nguyen, T. V. Nguyen, and M. T. Tran, ``MobileNet-SA: Lightweight CNN with self attention for sketch classification,'' in \textit{Pacific-Rim Symposium on Image and Video Technology (PSIVT) 2023}, Springer, 2023, pp.~115--129.

\bibitem{ref27}
L. Li, C. Zou, Y. Zheng, Q. Su, H. Fu, and C. L. Tai, ``Sketch-R2CNN: An RNN-rasterization-CNN architecture for vector sketch recognition,'' \textit{IEEE Transactions on Pattern Analysis and Machine Intelligence}, vol.~44, no.~9, pp.~4671--4685, 2020.

\bibitem{ref28}
Google Creative Lab, ``Quick, Draw! The Data,'' Google, 2017. [Online]. Available: https://quickdraw.withgoogle.com/data

\bibitem{ref29}
F. Liu, M. Lu, and W. Zhang, ``Comparative analysis of deep learning-based models in sketch recognition,'' in \textit{Proc. ICAIC 2024}, Atlantis Press, 2024.

\end{thebibliography}
